\section{Implementierung}

Implementiert ist ein Software-Prototyp für die reproduzierbare Datenaufbereitung, Detektion und Ergebnisdarstellung. Die physische Sortieranlage ist davon klar getrennt.

\subsection{Konfiguration und Kommandozeile}
Projektweite Klassen, Pfade sowie Standardwerte für Training und Inferenz liegen in \texttt{mira.yaml}. Eine modulare Kommandozeile stellt unter anderem Befehle zum Zusammenführen und Prüfen von Datensätzen, zum Trainieren, Exportieren und Evaluieren von Modellen sowie für Live-Inferenz und Dashboard bereit. Experimentparameter können in YAML-Dateien festgehalten und dadurch erneut ausgeführt werden.

\subsection{Datensatzaufbereitung}
Der generische Merger der Kommandozeile liest Quellen aus einem YAML-basierten Register ein. Er kann vorhandene YOLO-Daten übernehmen, Klassen-IDs abbilden und COCO-Annotationen in das YOLO-Format umwandeln. Beim Zusammenführen entstehen Trainings- und Validierungsordner sowie eine passende Dataset-YAML; ein Dry-Run und eine Strukturprüfung erlauben Kontrollen vor dem Training.

Der tatsächlich verwendete Bestand \texttt{merged\_mira\_balanced} wurde dagegen mit dem separaten Builder \texttt{scripts/build\_balanced\_dataset.py} erzeugt. Dieser lädt die vier festgelegten Quellen, wendet deren projektspezifische Klassenabbildungen und Splitregeln an, wählt die Trainingsbilder mit Seed~42 annähernd nach Boxanzahl aus, entfernt dateiidentische Bilder anhand von SHA-256-Prüfsummen und schreibt zusätzlich Testsplit und Herkunftsmanifest. Der Registry-Merger und dieser Builder sind somit zwei unterschiedliche Aufbereitungswege.

\subsection{YOLO-Inferenz und Tracking}
Die Live-Inferenz lädt Ultralytics-kompatible Detektoren und verarbeitet Bilder einer USB-Kamera. Ausgegeben werden Begrenzungsrahmen, Klassen und Konfidenzen. Konfidenz-, IoU- und Bildgrößenschwellen sind konfigurierbar; unsichere Treffer können in der Visualisierung gesondert markiert werden.

ByteTrack ist für PyTorch-Modelle (\texttt{.pt}) optional eingebunden und vergibt persistente Objekt-IDs über aufeinanderfolgende Frames. Für TFLite-Modelle ist Tracking im aktuellen Inferenzpfad deaktiviert. Damit ist ByteTrack Bestandteil des PT-Prototyps, aber kein nachgewiesener Teil einer Edge-Pipeline.

\subsection{Dashboard-Prototyp}
Das Control Center ist als FastAPI-Anwendung mit REST- und WebSocket-Schnittstellen umgesetzt. Es kann Kamerastatus, Modellwahl, Videoframes, Detektionen und Laufzeitmetriken an eine Browseroberfläche übertragen. Der aktuelle Stand ist ein Dashboard-Prototyp zur Beobachtung und Konfiguration der Inferenz, keine Steuerung einer Sortiermaschine.

\subsection{Modellexport}
Trainierte YOLO-Modelle können über die Pipeline nach FP32-TFLite, INT8-TFLite oder ONNX exportiert werden. Für den INT8-Export wird ein Kalibrierungsdatensatz angegeben. Die erzeugten Formate werden getrennt bewertet; aus einem erfolgreichen Export folgt noch keine Aussage zur Geschwindigkeit auf der Zielhardware.

\subsection{Nicht implementierte Systemteile}
Eine Ausführung auf Raspberry Pi, die Robotermechanik und eine serielle Verbindung zur Aktorik sind \textbf{nicht implementiert beziehungsweise ausstehend}. Entsprechend existieren derzeit weder ein geschlossener Regelkreis noch ein getestetes Handshake-Protokoll zwischen Erkennung und Roboter.
